\section*{الفصل \ref{ch:b2:liquid-vapour-diagrams} --- المخططات الثنائية سائل--بخار والتقطير}

\begin{solution}{exo:b2:liquid-vapour-diagrams:1}
البنزين \qty{80.1}{\degreeCelsius}، والتولوين \qty{110.6}{\degreeCelsius}. والسائل
ذو التركيب 0.2 يبدأ بالغليان عند نحو \qty{102}{\degreeCelsius}؛ وتركيب أول
بخار له 0.38.
\end{solution}

\begin{solution}{exo:b2:liquid-vapour-diagrams:2}
$p = 0.5 \times 1.010 + 0.5 \times 0.388 = \qty{0.699}{bar}$؛
$y = 0.505/0.699 = 0.722$.
\end{solution}

\begin{solution}{exo:b2:liquid-vapour-diagrams:3}
$n_V/n = (0.30 - 0.20)/(0.45 - 0.20) = 0.40$: \qty{4.0}{mol} من البخار 
و\qty{6.0}{mol} من السائل. التحقق: $6.0 \times 0.20 + 4.0 \times 0.45 = 3.0 = 10
\times 0.30$.
\end{solution}

\begin{solution}{exo:b2:liquid-vapour-diagrams:4}
(أ) $v = 2 + 2 - 1 = 3$: $T$ و$p$ و$x$ حرة. (ب) $v = 2 + 2 - 2 = 2$؛ وعند
ضغط ثابت 1. (ج) يضيف \omterm{def:b2:liquid-vapour-diagrams:azeotrope}{الأزيوتروب} العلاقة $x = y$: $v = 1$؛ وعند ضغط ثابت
0 --- فيغلي \omterm{def:b2:liquid-vapour-diagrams:azeotrope}{الأزيوتروب} عند درجة حرارة ثابتة، كالسائل النقي،
مع أن تركيبه يتغير مع الضغط.
\end{solution}

\begin{solution}{exo:b2:liquid-vapour-diagrams:5}
عند \omterm{def:b2:liquid-vapour-diagrams:bubble-dew}{نقطة الندى} $yp/p_1^* + (1 - y)p/p_2^* = 1$. عند \qty{100}{\degreeCelsius}:
$0.3 \times 1.013/1.80 + 0.7 \times 1.013/0.742 = 1.124 > 1$ (بارد جدًا)؛ وعند
\qty{105}{\degreeCelsius}: $0.148 + 0.824 = 0.971 < 1$. وبالاستيفاء، $T \approx
100 + 5 \times 0.124/0.153 = \qty{104}{\degreeCelsius}$. وتركيب أول قطرة $x =
yp/p_1^* \approx 0.3 \times 1.013/2.0 = 0.15$.
\end{solution}

\begin{solution}{exo:b2:liquid-vapour-diagrams:6}
لأول القطرات تركيب البخار فوق السائل المغلي،
$y = 0.71$، عند \qty{92.1}{\degreeCelsius}. ولأن البنزين يغادر بالأفضلية، يزداد
السائل فقرًا به، فترتفع درجة غليانه ويزداد القطير
فقرًا أيضًا؛ وآخر القطرات تولوين شبه نقي. فالتقطير البسيط
مرحلة توازن واحدة، تتكرر على سائل لا يكف عن التغير: فكل
قطفة خليط.
\end{solution}

\begin{solution}{exo:b2:liquid-vapour-diagrams:7}
التغذية 0.05، في جهة الماء من \omterm{def:b2:liquid-vapour-diagrams:azeotrope}{الأزيوتروب}: تعطي القمة \omterm{def:b2:liquid-vapour-diagrams:azeotrope}{الأزيوتروب}
(0.88 بالمولات، و95~\% كتلةً)، وتحتفظ المغلاة بالماء. والتغذية 0.95، في جهة
الإيثانول: تعطي القمة أيضًا \omterm{def:b2:liquid-vapour-diagrams:azeotrope}{الأزيوتروب}، وهو صاحب أدنى نقطة غليان،
وتحتفظ المغلاة بالإيثانول النقي.
\end{solution}

\begin{solution}{exo:b2:liquid-vapour-diagrams:8}
يغلي الخليط حين $0.10 + p^*_{\text{ماء}} = 1.013$، أي
$p^*_{\text{ماء}} = \qty{0.91}{bar}$، قرب \qty{97}{\degreeCelsius}. ونسبة الكتلتين
$0.10 \times 136/(0.91 \times 18.02) = 0.83$: أي \qty{0.83}{g} من الزيت العطري لكل غرام
من الماء.
\end{solution}

\begin{solution}{exo:b2:liquid-vapour-diagrams:9}
في درجة حرارة الغرفة، طبقتان L$_2$ وL$_1$ (فالتركيب 0.4 يقع داخل
الفجوة). وعند التسخين تغلي الطبقتان معًا عند درجة حرارة \omterm{def:b2:liquid-vapour-diagrams:miscibility-gap}{الأزيوتروب غير
المتجانس}، التي تبقى ثابتة؛ ولأول بخار التركيب
الأزيوتروبي (النقطة). ولأن 0.4 يقع في جهة L$_2$ من النقطة، تُستنفد الطبقة
L$_1$ أولًا؛ ثم ترتفع درجة الحرارة، ويزداد L$_2$ المتبقي
فقرًا بالمكوّن 1 على طول منحنى فقاعته ويتبع البخار \omterm{def:b2:liquid-vapour-diagrams:bubble-dew}{منحنى الندى}، حتى
تبلغ النقطة \omterm{def:b2:liquid-vapour-diagrams:bubble-dew}{منحنى الندى}: فتتبخر آخر قطرة.
\end{solution}

\begin{solution}{exo:b2:liquid-vapour-diagrams:10}
$N_{\min} = \ln(99 \times 99)/\ln 2.47 = 9.19/0.904 = 10.2$، مقابل 6.5: فالانتقال
من نقاوة 95~\% إلى 99~\% عند الطرفين يكلّف أكثر من نصف عدد الأطباق
مرة أخرى. فكل عامل إضافي على $x/(1 - x)$ يكلّف العدد نفسه من الأطباق،
والنقاوة تقاس بالشائبة المتبقية، التي لا تنخفض إلا هندسيًا.
\end{solution}

\begin{solution}{exo:b2:liquid-vapour-diagrams:11}
\omterm{def:b2:liquid-vapour-diagrams:binary-diagram}{المخطط المتساوي الضغط} وحده لا يعطيه: فلا بد من $p_1^*$ و$p_2^*$ عند
\qty{80}{\degreeCelsius}، من معادلتي أنطوان (1.010 
و\qty{0.388}{bar}). ثم $p = p_2^* + x(p_1^* - p_2^*)$ و$p = 1/(y/p_1^* + (1 -
y)/p_2^*)$. والسائل مستقر عند الضغط المرتفع (فضغط البخار
يكثفه)، وعند درجة الحرارة المنخفضة على \omterm{def:b2:liquid-vapour-diagrams:binary-diagram}{المخطط المتساوي الضغط}: فالمخططان
مقلوبان أحدهما بالنسبة إلى الآخر.
\end{solution}

\begin{solution}{exo:b2:liquid-vapour-diagrams:12}
لأن $y$ يتزايد مع $x$ و$y(1 - y) > 0$، فإن إشارة $dp/dx$ هي إشارة $y - x$:
فيرتفع الضغط مع $x$ حيث يكون البخار أغنى بالمكوّن 1. ومن أجل \omterm{def:b2:chemical-potential:ideal-mixture}{خليط
مثالي} $y > x$ لكل $0 < x < 1$ (\cref{prop:b2:liquid-vapour-diagrams:richer}):
فيكون $p$ رتيبًا تمامًا، بلا قيمة قصوى، ولا \omterm{def:b2:liquid-vapour-diagrams:azeotrope}{أزيوتروب}.
\end{solution}

\begin{solution}{pb:b2:liquid-vapour-diagrams:1}
\textbf{1.} $T = B/(A - \log_{10}1.013) - C$: البنزين $1203.835/4.01253 + 53.226 =
\qty{353.2}{K}$، أي \qty{80.1}{\degreeCelsius}؛ والتولوين $1343.943/4.07266 + 53.773 =
\qty{383.8}{K}$، أي \qty{110.6}{\degreeCelsius}.
\textbf{2.} عند \qty{363.15}{K}: $p_1^* = \qty{1.361}{bar}$، $p_2^* = \qty{0.542}{bar}$.
\textbf{3.} $x = (1.013 - 0.542)/(1.361 - 0.542) = 0.575$؛ $y = 0.575 \times
1.361/1.013 = 0.773$.
\textbf{4.} $x = (1.013 - 0.742)/(1.800 - 0.742) = 0.256$؛ $y = 0.256 \times
1.800/1.013 = 0.455$.
\textbf{5.} نقاط الفقاعة (0، 110.6)، (0.256، 100)، (0.575، 90)، (1، 80.1)؛ ونقاط
الندى (0، 110.6)، (0.455، 100)، (0.773، 90)، (1، 80.1) --- وهي العدسة التي في
الشكل.
\textbf{6.} $\frac12(p_1^* + p_2^*) = 1.013$: عند \qty{90}{\degreeCelsius} 0.952، وعند
\qty{95}{\degreeCelsius} 1.102؛ وبالاستيفاء، $90 + 5 \times 0.061/0.150 \approx
\qty{92}{\degreeCelsius}$ (\qty{92.1}{\degreeCelsius} بالضبط).
\textbf{7.} $v = 2$، و1 عند ضغط ثابت: فكل درجة حرارة تحدد $x$ و$y$،
ولهذا يتكوّن المخطط من منحنيين.
\textbf{8.} $x = (1.013 - 0.636)/(1.569 - 0.636) = 0.404$؛ $y = 0.404 \times
1.569/1.013 = 0.626$.
\textbf{9.} $n_V = 100 \times (0.500 - 0.404)/(0.626 - 0.404) = \qty{43}{mol}$،
$n_L = \qty{57}{mol}$.
\textbf{10.} البخار $43 \times 0.626 = \qty{26.9}{mol}$؛ والسائل $57 \times 0.404
= \qty{23.0}{mol}$؛ والمجموع \qty{49.9}{mol} $\approx 50$.
\textbf{11.} $26.9/50 = 54~\%$ من البنزين، في بخار نسبته 63~\% فقط.
\textbf{12.} عند \qty{100}{\degreeCelsius} يكون للبخار المتوازن $y =
0.455 < 0.5$: فنقطة التغذية تقع فوق \omterm{def:b2:liquid-vapour-diagrams:bubble-dew}{منحنى الندى}، والتغذية كلها
بخار ولا يُفصل شيء.
\textbf{13.} عند نقطة ندى التغذية، \qty{98.8}{\degreeCelsius} (بطريقة
\cref{exo:b2:liquid-vapour-diagrams:5}).
\textbf{14.} $\alpha = p_1^*/p_2^*$؛ عند \qty{80.1}{\degreeCelsius}، $1.013/0.390 =
2.60$.
\textbf{15.} عند \qty{110.6}{\degreeCelsius}، $2.377/1.013 = 2.35$.
\textbf{16.} $\sqrt{2.60 \times 2.35} = 2.47$.
\textbf{17.} $y = xp_1^*/p$ و$1 - y = (1 - x)p_2^*/p$؛ نقسم.
\textbf{18.} لا يُسحب شيء: فبين طبقين يحمل البخار الصاعد
والسائل النازل الكمية نفسها من المادة ومن كل مكوّن،
$V y_k = L x_{k+1}$ مع $V = L$، ومنه $x_{k+1} = y_k$.
\textbf{19.} مع $r = x/(1 - x)$: $r(y_k) = \alpha\,r(x_k)$ و$x_{k+1} = y_k$، ومنه
$r(x_{k+1}) = \alpha\,r(x_k)$ و$r(x_D) = \alpha^N r(x_B)$؛ $N = \ln[r(x_D)/
r(x_B)]/\ln\alpha$.
\textbf{20.} $N_{\min} = \ln(19 \times 19)/\ln 2.47 = 5.889/0.904 = 6.5$.
\textbf{21.} يجب أن يعطي العمود ناتجًا، فيعمل بنسبة ارتداد منتهية،
تتطلب أطباقًا أكثر؛ والطبق الحقيقي لا يبلغ التوازن (كفاءته
أقل من واحد).
\textbf{22.} سبع درجات، تنتهي آخرها عند 0.034، تحت 0.05: أي 6.5 طبقًا في
عدّ متصل، وسبعة أطباق كاملة (والمغلاة تُحسب واحدًا منها).
\textbf{23.} $x = (95/46.07)/(95/46.07 + 5/18.02) = 0.881$.
\textbf{24.} في القمة، \omterm{def:b2:liquid-vapour-diagrams:azeotrope}{الأزيوتروب} في أحسن الأحوال (0.88 بالمولات)؛ وفي القاع،
الماء. فخمسون طبقًا لا تعبر \omterm{def:b2:liquid-vapour-diagrams:azeotrope}{الأزيوتروب}.
\textbf{25.} عند الارتداد الكلي، $\boldsymbol{N_{\min} \approx 6.5}$ طبقًا
نظريًا.
\end{solution}
